\subsection{A dyadic multiplier test for the congruence-recognition problem}

\begin{statement}{A.54}{Proved}{A general recognition criterion with one coordinate supported at two}
Let $m\ge5$ be odd, and let
\[
 \lambda_0=2^m\quad\hbox{or}\quad\lambda_0=1+2^m.
\]
Write $M=2^m-1$ in the first case and $M=2^m+1$ in the second.
Let $s$ be a noncuspidal Selmer section whose full integral
Legendre representation is conjugate to that of the point
$\lambda_0$ by an element of the geometric closure $H$ in A.17.
Every prime $p\mid M$ at which $s_p$ is cuspidal satisfies
\[
             p\nmid(m-4),\qquad
             \left(\frac{4-m}{p}\right)=1.
\tag{A.54.1}
\]
Consequently, when no prime divisor of $M$ satisfies (A.54.1),
$s$ is precisely the point section $s_{\lambda_0}$.

In particular, the conclusion holds for $\lambda_0=33$:
$M=33=3\cdot11$ and $4-m=-1$ is a nonsquare modulo both
prime divisors.  This example settles its exact A7 recognition
question.  The proof uses the determinant-one hypothesis on the
conjugator; it does not assert this conclusion for every
unrestricted integral $\operatorname{GL}_2$-intertwiner.
\end{statement}
\begin{proof}
We use multiplicative notation for completed Kummer classes and
additive notation for their valuations and logarithms.  The
integral comparison with $E_{\lambda_0}$ gives finite ramification
as in A.14 and preserves the two ordered nonzero points
$P_1=(0,0),P_2=(1,0)$ of order two.  In particular A.16 gives
\[
        A_s\equiv\lambda_0,\qquad
        B_s\equiv1-\lambda_0
                 \pmod{(\widehat{\Q^\times})^2}.
\tag{A.54.2}
\]
Define the coordinate $C$ and the cusp $c$ by
\[
\begin{array}{c|c|c|c|c}
\lambda_0&C&c&\text{other coordinate}&\text{opposite cusp }d\\ \hline
2^m&A_s&0&B_s&1\\
1+2^m&B_s&1&A_s&0.
\end{array}
\tag{A.54.3}
\]
Thus $C$ vanishes at $c$ and has value one at $d$.

For an odd prime dividing $M$, the local curve
$E_{\lambda_0}$ has multiplicative reduction, with the collision
of roots belonging to the cusp $d$.  A.41 identifies this label
using the canonical saturated cyclotomic line in $T_2$ and its
reduction modulo two.  It follows that $s_p$ is either an ordinary
point with that contact or a section in the packet of $d$.
In both cases the valuation of $C$ at $p$ is zero.  At every
other odd prime the realized curve has good reduction; the
local comparison excludes a cuspidal section or a nonintegral
ordinary point there, by A.37 and the local monodromy calculation
in A.14.  The valuation of $C$ is again zero.

At two, (A.54.2) says that the valuation of $C$ is odd in
$\widehat\Z/2$, and that of the other coordinate is even.
For an ordinary point $x\in U(\Q_2)$ the valuation pair is
exactly one of
\[
      (a,0)\ (a>0),\qquad (0,b)\ (b>0),\qquad
      (a,a)\ (a<0),\qquad (0,0).
\tag{A.54.4}
\]
This follows directly from $x+(1-x)=1$ and the nonarchimedean
valuation.  The parity in (A.54.2) leaves only the positive
valuation in the $C$ coordinate and zero in the other coordinate.
For a cusp section, the valuation at cusp zero has second
coordinate zero, the valuation at cusp one has first coordinate
zero, and at infinity the two valuations are equal.  The same
parity therefore leaves only cusp $c$, with the other valuation
zero.  This proves, in particular, that for $\lambda_0=33$ a
cuspidal localization at two must have label one; neither the
label nor this assertion uses an ordinary-point assumption.

The completed rational multiplicative group has its valuation
description by the sign and restricted families in
$\widehat\Z$.  Since the valuation of $C$ is zero at every odd
prime, we have exactly
\[
                    C=\varepsilon\,2^\beta,
       \qquad \beta\in\widehat\Z,
\tag{A.54.5}
\]
where $\varepsilon=1$ in the first row of (A.54.3) and
$\varepsilon=-1$ in the second, as fixed by (A.54.2).
The valuation of the other coordinate at two is zero, as
just proved.  Thus (A.54.5) is a global identity of completed
classes, not merely a description of one local component.

Suppose that $s_p$ is cuspidal for an odd prime $p\mid M$.
Its label is $d$, so the unit $C$, regular at $d$ with value
one, has trivial localization there.  Apply the coefficient-$p$
Kummer map and then the $p$-adic logarithm to (A.54.5).  This gives
\[
                      \beta_p\log_p(2)=0.
\tag{A.54.6}
\]
Here the logarithm kills the sign.  Also $\log_p(2)\ne0$:
the logarithm is injective on $1+p\Z_p$, and a unit with zero
logarithm is a root of unity.  The positive rational number two
is not a root of unity, since a relation $2^r=1$ in $\Q_p$
would be the same false equality of rational numbers.
Therefore
\[
                              \beta_p=0.
\tag{A.54.7}
\]

The localization at two cannot now be ordinary.  Indeed its
valuation of $C$ would be an ordinary positive integer $n$, so
the valuation homomorphism at two applied to (A.54.5) would
give $\beta=n$ in $\widehat\Z$.  Its $p$-component would be
nonzero, contrary to (A.54.7).  Thus $s_2$ has label $c$.

We compare this cusp localization with the realized elliptic
curve at two, keeping the integral pairing.  The Tate parameter
of the Legendre family at $c$, after its constant quadratic
splitting twist when $c=0$, has expansion
\[
                  q=\frac{C^2}{256}\,F(C),
                       \qquad F(0)=1.
\tag{A.54.8}
\]
This is the formal inversion of
$j(q)=q^{-1}+744+\cdots$ using the displayed Legendre $j$
function; the leading coefficient is $1/256$ in either row.
The factor $F$ is a formal unit with compatible roots having
constant term one.  On the cusp decomposition group its
Kummer class is therefore zero.  The exact coefficient-$p$
Tate extension class of the cusp section is consequently
\[
                     2\kappa_p(C)-8\kappa_p(2).
\tag{A.54.9}
\]
This calculation retains the constant 256; discarding it would
lose the present obstruction.

For the actual rational parameter, $\lambda_0^2-\lambda_0+1$
is odd, and the two coordinate valuations are $m$ and zero.
It follows that
\[
             v_2(j(E_{\lambda_0}))=8-2m<0,
                 \qquad v_2(q_0)=2m-8.
\tag{A.54.10}
\]
Choose a common finite extension $F/\Q_2$ over which the two
potentially multiplicative representations are split Tate
representations, as in A.34.  Let $e$ be its ramification index.
After restriction to $G_F$, the valuation of (A.54.9) is
$e(2\beta_p-8)=-8e$, while that of the actual Tate parameter
is $e(2m-8)$.

The chosen global conjugator lies in $H$, so its coefficient-$p$
component has determinant one.  Choose the usual Tate bases
whose Weil pairings have value the chosen compatible roots
of unity; the determinant remains one in these bases.
The nonzero Tate extension has a unique saturated cyclotomic
line.  The intertwiner must preserve it and hence has the form
\[
                        g_p=\begin{pmatrix}a&b\\0&d\end{pmatrix},
                 \qquad a,d\in\Z_p^\times,\quad ad=1.
\tag{A.54.11}
\]
For clarity, if the two Tate matrices are
$\left(\begin{smallmatrix}\chi&c_0\\0&1\end{smallmatrix}\right)$
and $\left(\begin{smallmatrix}\chi&c_s\\0&1\end{smallmatrix}\right)$,
the equality $g_p\rho_0=\rho_s g_p$ says
\[
                     d c_s=a c_0+(1-\chi)b.
\]
Passing to cohomology and then to the valuation gives
\[
       \frac{-8e}{e(2m-8)}
          =\frac{-4}{m-4}=a/d=a^2
                         \quad\hbox{in }\Q_p^\times.
\tag{A.54.12}
\]
In particular this ratio is a square unit.  Since $p$ is odd,
this is equivalent to (A.54.1): it first requires
$p\nmid(m-4)$, and then requires $4-m$ to be a square modulo
$p$.  This proves the asserted restriction on every odd
cuspidal prime.

If no prime divisor of $M$ passes (A.54.1), there are no
cuspidal localizations at odd primes.  The preceding valuation
argument leaves only label $c$ at two.  The one-label recognition
result A.37 therefore makes $s$ a rational point section.

We finally check that this point has parameter $\lambda_0$,
rather than just the same integral Tate system.  Write it as
$s_\lambda$.  Faltings's isogeny theorem gives a rational
isogeny between $E_\lambda$ and $E_{\lambda_0}$.
For an elliptic curve over $\Q$ its rational endomorphism
ring is $\Z$: an endomorphism acts on its rational invariant
differential by a rational scalar, and its complex-analytic
scalar is an algebraic integer, so the scalar is an integer.
Thus the group of rational homomorphisms between these two
isogenous curves is free of rank one.  Choose its primitive
isogeny generator $\phi$.

At each coefficient prime $\ell$, Schur's lemma and A.12
express $T_\ell\phi$ as a scalar times the given integral
intertwiner.  Its two elementary divisors therefore have the
same $\ell$-adic valuation.  If that common valuation were
positive, $\phi$ would kill $E_\lambda[\ell]$ and factor over
$\Q$ through $[\ell]$, contradicting primitivity.  Every
common valuation is zero, so $\deg\phi=1$.  Moreover its
reduction modulo two has the specified labels: its scalar
relative to the integral intertwiner is a $2$-adic unit,
which reduces to one in $\mathbf F_2$.
An isomorphism between the two Legendre equations preserving
the origin has $x$ coordinate $x\mapsto u^2x+r$.
Preservation of $(0,0)$ forces $r=0$, and preservation of
$(1,0)$ forces $u^2=1$.  Preservation of the third nonzero
point of order two then gives $\lambda=\lambda_0$.
This proves the exact point-section conclusion.
\end{proof}

\noindent\emph{Scope of the multiplier computation.}
With a general labeled integral intertwiner, (A.54.11) has
$ad=\delta_p$ and (A.54.12) reads
$-4/(m-4)=a^2/\delta_p$.  It then computes a necessary
multiplier squareclass, rather than forcing (A.54.1).
This is the precise place where the geometric-conjugacy
hypothesis in A7 enters the argument.


\noindent\emph{Exact finite certificate.}
\texttt{code/legendre\_dyadic\_certificate.py} computes
$j(E_{33})=1180932193/4356$, its valuation $-2$ at two,
the ratio $-4$, and its exclusion from unit squares modulo
$3,9,11,121$. It also verifies the intertwiner matrix identity
and the coefficient $1/256$ at both formal cusps.
The actual output is retained alongside the program.
